\subsection{Evaluation of the Inundation Model}
The Inundation Model has ben a good tool for simulating flood propagation at certain water levels and has been reasonable for simulating a previous storm surge event that occurred in Denmark in October 2023. There are a number of advantages and disadvantages of the mode, which will be evaluated and discussed.
The biggest advantage the model possesses is its ease of use and relatively simple structure. It is also advantageous that the model operates with a graphical user interface, which means that it does not require knowledge of programming or using a command line interface, which some of the more advanced non-GIS models utilise. This makes it easy for individuals, students, and organisations to use the model without requiring broader computer knowledge. In addition, it is an advantage that the model operates inside of a GIS-environment and can thus be combined with other geospatial data, such as as land use data, and other point, line and polygon data to investigate flooding risks on sectors and critical infrastructure. 
The model is also relatively fast in processing, if the hardware is good. The processing duration for this study has been a total of 16 hours and 45 minutes and has therefore been considerably longer than if used on more powerful hardware. It is expected to twice as fast on more powerful hardware.\\

The model also is comprised of several tools that have not been utilised in this study, including a tool that can determine water entry points. In combination with this tool, the model can be rerun to determine whether mitigation procedures have been successful. Another tool is one that is able to calculate water depths at a given location and be used to perform damage cost analysis. Coupled with the relatively fast simulation time, this is a fast working approach that can complement the municipalities' storm surge and coastal planning. These two additional tools contribute to the model's usefulness in planning proactive storm surge protection measures.\\

Despite this, there are also a number of limitations to the model. The model is a 100\% static numerical model and this therefore means that there is no possibility of incorporating a temporal scale. Many storm surges occur over a period of a few hours to a couple days. This means that it takes time for a flooding to reach its full extent and often mitigation measure will have been implemented to slow down the spread of water, as seen in Aabenraa. The model therefore provides an idea of what a flood will look like, given an infinitely long time period and without intervention from society. The missing temporal scale also means that model does not take correct hydrodynamic behaviour into account. 

Furthermore, the model doesn't take other climatic aspects that happend during a storm surge into account. This includes effects of wind, wave contributions and rainfall. The model can therefore only provide an overview of the flooding based on the water level. Thus, if there is a protective dike along a stretch of coast that is higher than the simulated water level, it will not be flooded by the model, whereas in a real storm surge event, it is likely that waves and wind will push water over the dike and flood the area behind it. In the same vein, the model is not suitable for handling compound flooding, where an area is flooded from precipitation and the sea at the same time, since the Inundation uses a singular source for its calculation.
These are elements that the more advanced models such as Deltares SFINCS model and DHI's MIKE21/3 models are able to simulate, as they include many climatological parameters in their calculations.\\

The model is also highly dependent on the availability of high-resolution data. The digital terrain model (DTM) should preferably be in high resolution with the smallest cell size possible \citep{seenath_effects_2018, williams_geographic_2022} to be more realistic and reduce accuracy uncertainty. In Denmark this is not a problem, as a high-quality DTM can be downloaded for free from the Danish Agency for Climate Data, but in other countries where a good terrain model has yet to be established this can present a sizeable issue. A DTM with a cell size of more than five metres for local scales and more than 50 metres on regional scales will imply uncertainty about the accuracy of the simulated flooding \citep{williams_geographic_2022}.\\ 
Comprehensive and updated datasets of hydrological adaptations are also a requirement for the model to simulate correct hydrological behaviour \citep{bales_sources_2009}. The hydrological adaptations are less important than the quality of the DTM if they do not exists for the area, as it is possible to collect this data in the field and manually digitise the adaptations.\\

Finally, it is worth mentioning that the model is currently exclusive to the ArcGIS Pro software. This means that the model works very well in that specific environment and can easily be integrated into other parts of Esri's ArcGIS environment, such as ArcGIS Online, Enterprise and Field Maps, for working in the field and across collaborators. Being limited to ArcGIS also means that the model is locked behind a steep paywall, especially for commercial use, and is outside the potential use of Danish municipalities and other private individuals, who cannot meet the high price. A solution to this would be to develop or translate the model to another GIS software, such as QGIS, a free and open-source software (FOSS) that many Danish municipalities use for spatial planning and management.\\

In summary, the Inundation Model is a model with a lot of potential to be utilised in initial storm surge analysis. The ease of use meant that the model cold be easily distributed to many people of interest in municipalities, companies and private individuals. Especially if the model can be transferred to a FOSS distribution or a non-commercial GIS product. The model should be used in combination with more advanced models and tools, but is not able to replace the need for more complex hydrodynamic models due to the Inundation Model's static approach.